\section{Possible paper directions, ranked by how well the evidence supports them}\label{sec:directions}

Each direction lists the claim, the evidence that already exists (with its level), what is missing, and an estimate of the cost. Cost is given in episodes; on this laptop one evaluation episode takes about 13 seconds (the 400-episode check of Section~\ref{sec:weather} took 89 minutes), so 200 episodes cost about 45 minutes.
The directions are not exclusive: A, B and C together form one coherent paper (the ``thorough evaluation'' option at the end).

\subsection{Direction A: pitfalls of closed-loop evaluation of learned driving policies (strongest support)}
\textbf{Claim.} Common practices in evaluating driving policies in simulation (small test sets, validation-based selection, unpaired comparisons, unvalidated sensor perturbations, non-deterministic traffic) produce wrong conclusions, with a concrete catalogue of cases and a protocol that avoids them.
\textbf{Evidence in hand.} (i)~A model's success of 0.975 on 40 episodes became 0.895 on 200 (selection of the best of several candidates; reported). (ii)~Reinforcement-learning validation of 1.0 at every checkpoint while the 200-episode test shows a significant loss of 0.065--0.115 (reported). (iii)~An invalid perturbation measurement from a stale sensor mount, detected only by a ``identical across conditions'' check (documented). (iv)~The same model and seed giving different episodes in different process histories (diagnostic).
(v)~Ties in 10-episode validation resolved to round 0 in four runs (documented). (vi)~Paired exact tests detect a 0.075 difference at $n=200$ that overlapping Wilson intervals do not.
\textbf{Missing.} A quantitative study of (i) by resampling from the per-episode data that already exists (selection bias as a function of the number of candidates and $n$); a determinism fix or a documented seeding protocol; a short checklist. \textbf{Cost:} mostly analysis of existing data; no new episodes needed for the first version.
\textbf{Risk:} a methodology paper is judged on generality; the cases come from one simulator and one policy family.

\subsection{Direction B: where does perception degradation start to matter? A measured dose-response study (strong support, cheap)}
\textbf{Claim.} Under a measured model of detector and stereo-depth degradation under fog and rain, sparse-traffic driving is unaffected, and the paper maps the region of (traffic density, vehicle detection retention) in which success begins to fall, i.e.\ the conditions under which any weather-adaptive method would have something to adapt to.
\textbf{Evidence in hand.} The measurement (Figure~\ref{fig:perception}, pilot level: 72--144 vehicle-frames per condition) with its finding that loss is detection, not range error; the calibrated environment; the null result at $n=200$ for four models in four weather cells and at $n=100$ against clear (pre-registered, reported); the dropout proxy (pilot) with the same null.
My literature search found no 2025--2026 arXiv paper that evaluates weather-degraded perception in MetaDrive (Section~\ref{sec:related}), so the combination is open; this is a statement about my search, not proof of novelty.
\textbf{Missing.} A grid of retention $\in\{1,0.75,0.5,0.25,0\}$ times density $\in\{0.05,0.1,0.2\}$ for two or three models; a larger measurement sample; ideally the same measurement with a second detector. \textbf{Cost:} 25 cells $\times$ 100 episodes $\approx$ 9 hours per model.

\subsection{Direction C: the teacher--student information gap in bird's-eye-view driving (strong potential, one decisive experiment)}
\textbf{Claim.} A large part of the remaining failures of an imitation-trained bird's-eye-view policy comes from what its grid cannot show (no heading, 17.5\,m window, centre rule), not from imitation error; a heading-aware, wider grid or a teacher restricted to the student's information closes the gap.
\textbf{Evidence in hand.} The headline model fails 15 of 200; the collision-aware teacher solves 14 of them; seven of the failures are collisions in the first 10 seconds in which every other model also collides (reported); the fidelity replay of two crashes (diagnostic).
Related 2025--2026 work frames exactly this issue (learner--expert asymmetry, unrealizable teachers; Section~\ref{sec:related}).
\textbf{Missing (decisive).} (1)~Evaluate the existing student of the collision-aware teacher on the 200-episode block (45 minutes); (2)~a teacher that sees only the student's window (to separate information gap from imitation error); (3)~a heading-aware grid (policy-side wrapper around the teammates' grid) and retraining; (4)~three training seeds. \textbf{Cost:} (1) 200 episodes, (2)--(3) about 1--2 days of training and 3 $\times$ 200 episodes.
\textbf{Risk:} the improvement may be small if the remaining failures are dominated by other effects; the experiment answers that either way.

\subsection{Direction D: a thorough evaluation of imitation versus reinforcement learning in a camera-derived bird's-eye-view setting}
\textbf{Claim.} In sparse and dense MetaDrive traffic, dataset aggregation from a privileged rule-based teacher beats few-demonstration behaviour cloning and beats reinforcement-learning fine-tuning, with a full account of why four reinforcement-learning attempts failed.
\textbf{Evidence in hand.} The 200-episode comparison of six models; the four reinforcement-learning attempts with their diagnoses; the reward-design audit trail.
\textbf{Missing.} (1)~Several training seeds per method (all results are single training seeds, so between-seed variance is unknown); (2)~a reinforcement-learning baseline from scratch with an algorithm better matched to MetaDrive (a third-party summary of the MetaDrive paper states that value-based methods suited driving better than policy-gradient ones; not checked in the full text); (3)~a stronger reinforcement-learning attempt with the repairs listed in Section~\ref{sec:rl}; (4)~dense-traffic numbers at $n\ge100$.
\textbf{Cost:} 3 seeds $\times$ about 2 hours of training per model, plus about 45 minutes per 200 evaluation episodes.

\subsection{Direction E: robustness of map-overlay policies to localisation error through redundancy}
\textbf{Claim.} A policy that reads lane geometry from a map overlay is brittle to an error in the car's pose against the map, but training with pose error \emph{and a second independent lane cue} removes the brittleness; augmentation alone does not.
\textbf{Evidence in hand.} Success 0.925 $\to$ 0.275 from 0 to 2\,m of error for the headline model; $0.88$--$0.95$ for the two-cue models over the same range; no benefit from augmentation alone (all pilot, 40 episodes per point, single training seed).
\textbf{Missing.} 200 episodes per point; replace the simulated second cue by the output of the real lane/stereo pipeline; a time-varying (not constant) pose error; multiple seeds. Related benchmark work perturbs GPS and odometry in closed loop (Bench2Drive-Robust, Section~\ref{sec:related}).
\textbf{Cost:} 5 levels $\times$ 200 episodes $\times$ 3 models $\approx$ 11 hours.

\subsection{Direction F: reliability-coupled caution as a speed--energy trade-off}
\textbf{Claim.} A single network conditioned on a perception-reliability value spans a controllable trade-off between speed, estimated energy and timeouts at roughly constant success.
\textbf{Evidence in hand.} Pilot at 40 episodes: speed 8.0 $\to$ 5.6--6.2\,m/s, energy 84.7 $\to$ 56--59\,Wh/km (30\% lower), success 0.85--0.925, timeouts 7.5--15\%.
\textbf{Missing.} $n=200$; an energy model validated against a published vehicle model; a longer horizon so that timeouts do not confound slow driving; a Pareto front over several caution values; a control that simply lowers the target speed of the stock driver model (to show that learning adds anything). \textbf{Cost:} about 9 hours. \textbf{Risk:} the strongest alternative explanation, ``a slower car uses less energy'', must be excluded by the speed-matched control.

\subsection{Direction G: a reproducible benchmark and dataset}
\textbf{Claim.} A resource: the evaluation protocol, seed blocks, per-episode outcomes of more than ten policies, the perception measurement and the failure catalogue, packaged for others.
\textbf{Evidence in hand.} Everything is stored (per-episode JSON, model manifest with hashes, scripts); the repository is organised for it (Section~\ref{sec:repomap}).
\textbf{Missing.} Push to the remote, licence check of the third-party weights, determinism fix (Section~\ref{sec:repro}), documentation of the environment. \textbf{Cost:} one to two days of engineering.

\subsection{Recommended combination}
The most defensible \emph{single} paper with the data that exists is a \textbf{thorough evaluation paper} built from A, B and the existing part of D, with C as its central analysis: \emph{``Imitation, reinforcement learning and perception degradation for bird's-eye-view driving in MetaDrive: a paired evaluation, a measured weather model, and the information gap that limits the students''}.
Its claims would be: (1) dataset aggregation from a privileged teacher beats few-demonstration cloning and reinforcement-learning fine-tuning (reported); (2) measured fog and rain remove detections but do not hurt sparse-traffic driving (reported, with the dose-response map from B added); (3) most remaining failures are information-gap failures (reported overlap, plus the experiment of C); (4) evaluation pitfalls (A).
The experiments to add before writing are, in priority order: (i)~the collision-aware student and the unmodified-grid student on the 200-episode block (about 1.5 hours); (ii)~two more training seeds for the headline model; (iii)~the dose-response grid; (iv)~the heading-aware grid test.
